\clearpage\thispagestyle{empty}
\pdfbookmark[0]{Quellen zu Komet (3. Juli 1815)}{comet1815jul-sources}
\begin{center}\textbf{Quellen zu Komet (3. Juli 1815)}\end{center}
\noindent Diese KI-Übertragung ins heutige Deutsch folgt dem Bericht mit der Überschrift Comet in den \href{https://gdz.sub.uni-goettingen.de/id/PPN236007467}{\textit{Werken}, Band~VI} von Gauß (Göttingen, 1874), S.~385--387. Der Erstdruck erschien unter Göttingen in den \textit{Göttingischen gelehrten Anzeigen}, Stück~105 (3. Juli 1815). Die Werkausgabe nennt S.~1041..1043. Im Exemplar der Bayerischen Staatsbibliothek reicht der vollständige Erstdruck von S.~\href{https://api.digitale-sammlungen.de/iiif/image/v2/bsb10538687_00367/full/1600,/0/default.jpg}{1041} bis S.~\href{https://api.digitale-sammlungen.de/iiif/image/v2/bsb10538687_00371/full/1600,/0/default.jpg}{1045}.
\par\smallskip\noindent Der Erstdruck gibt den aufsteigenden Knoten der parabolischen Bahn als $82\,43\,6$ an. Die Werkausgabe setzt Einheitenzeichen, $82^\circ\,43'\,6''$. Für die Lichtstärke am 20.~Juni steht im Erstdruck auf S.~\href{https://api.digitale-sammlungen.de/iiif/image/v2/bsb10538687_00370/full/1600,/0/default.jpg}{1044} $\frac{9}{10}$, in der Werkausgabe $\frac{3}{10}$. Der Bruchteil der Umlaufzeit beträgt im Erstdruck $\frac{1}{3}$, wobei die vorangehenden Ziffern blass sind. In der Werkausgabe beträgt die Umlaufzeit $77\frac{1}{8}$ Jahre. Der Text folgt den Lesarten der Werkausgabe. KI-Übertragung und KI-Satz.
\clearpage
